% Nouveau EGA V (cote 25): archive cote 25, site page 128 (PDF page 129).
% typescript. Draft compared: none.
% AI-transcribed by Claude Opus 5.5 (claude-opus-5-5, adaptive thinking, effort high), Message Batches API, batch ega25b (msgbatch_01A7n5YuXubUPNK1zsepi8RG), 2026-11-13; not yet checked by eye.
\scanpage{25/128}
\mspage{\hand{31}}

par une section inversible de $\underline{O}_S$ , donc $\hand{\phi}$ et $\hand{\psi}$ définissent le même point
de $\overset{\hand{\vee}}{P}$ à valeurs dans S. Comme tout point de $\overset{\hand{\vee}}{P}$ à valeurs dans S est défini,
localement sur S, par une section $\hand{\phi}$ de $\underline{E}_S$ qui ne s'annule pas (cf Chap II)
a) s'ensuit.   Pour prouver b), on note le

\uline{Lemme}   Soit D un syst lin de diviseurs sur X non vide loc de type fini
sur k alg clos, paramétré par $\overset{\hand{\vee}}{P}$(E), et considérons le morphisme correspon-%
dant  f: X-Z $\hand{\longrightarrow}$ P , où Z est l'ensemble des points fixes. Alors \var{\xout{xx}}{}{typewriter strike-out}{ega25_25_128-v0001} si
r = rang$_k$E \hand{$>0$,} il existe  r+1 point\hand{s}\var{\inserted[typed]{$x_i$}}{$x_i$}{insertion}{ega25_25_128-v0002}\var{\inserted{$(1\leq i\leq r+1)$}}{$(1\leq i\leq r+1)$}{insertion}{ega25_25_128-v0003} de \var{\xout{(Xx}}{}{typewriter strike-out}{ega25_25_128-v0004} X(k)-Z(k)  tels que les
f(x$_i$) forment une  ``base projective'' de P, i.e. tels que  pour toute
partie J de [1,r+1] ayant \hand{r} \var{\xout{mi}}{}{typewriter strike-out}{ega25_25_128-v0005} éléments,  les  f(x$_i$),i\hand{$\in J$}, \var{\xout{xxxx}}{}{typewriter strike-out}{ega25_25_128-v0006} ne sont pas
contenus dans un hyperplan de P.

On peut évidemment supposer Z=\hand{$\emptyset$} , \var{\deleted{X intègre (en le remplaçant
par une\var{\xout{Xxxxxx}}{}{typewriter strike-out}{ega25_25_128-v0008} composanteirréductible T de X muni de la structure induite
réduite).}}{}{deletion}{ega25_25_128-v0007} \var{\xout{OxxaxxxxxapplicationxxinjectivexxxE}}{}{typewriter strike-out}{ega25_25_128-v0009} Comme \var{\inserted[typed]{par \var{\over{4}{\unsure{3|5}}}{\unsure{3|5}}{writing over}{ega25_25_128-v0011}.1.}}{par \var{\over{4}{\unsure{3|5}}}{\unsure{3|5}}{writing over}{ega25_25_128-v0011}.1.}{insertion}{ega25_25_128-v0010} f(X) n'est contenu dans
aucun hyperplan de P, on en conclut d'abord l'existence de r points
$x_i$\var{\inserted{$(1\leq i\leq r)$}}{$(1\leq i\leq r)$}{insertion}{ega25_25_128-v0012} de X \var{\xout{telsxque}}{}{typewriter strike-out}{ega25_25_128-v0013}\var{\inserted[typed]{\unsure{teks|tels} que les f($x_i$) soient}}{\unsure{teks|tels} que les f($x_i$) soient}{insertion}{ega25_25_128-v0014} projectivement indépendants dans P, i.e. soient définies
par des formes linéairement indépendantes sur E.   Reste à prouver qu'il
existe un  $x_{r+1} = x$\var{\deleted{\illegible{1}}}{}{deletion}{ega25_25_128-v0015}   dans X(k) tel que  f(x) ne soit dans aucun des r
\unsure{hyperpbans|hyperplans} \var{\inserted[typed]{$H_i$}}{$H_i$}{insertion}{ega25_25_128-v0016} définis par les \var{\deleted{\unsure{f($x_i$)}}}{}{deletion}{ega25_25_128-v0017} systèmes de r-1 parmi les f(x$_i$). Or dans
le cas contraire, compte tenu du sorite Jacobson,  on aurait
f(X) $\hand{\subset\bigcup}_{\hand{i}}$ H$_i$ , donc \var{\xout{Xxseraitxlaxréunionxdesxxxx}}{}{typewriter strike-out}{ega25_25_128-v0018}\var{\xout{xxx}}{}{typewriter strike-out}{ega25_25_128-v0019} si X$_0$  est une composante
irréductible de X, f(X$_0$) serait contenu dans un des H$_i$ , ce qui contredit
\var{\over{\unsure{4}}{\unsure{3|5}}}{\unsure{3|5}}{writing over}{ega25_25_128-v0020}.1. - Ceci posé,  pour prouver b), on peut  évidemment supposer Z=\hand{$\emptyset$} ,
et utilisant \var{\over{\unsure{4}}{\unsure{3|5}}}{\unsure{3|5}}{writing over}{ega25_25_128-v0021}.1.  on est ramené à prouver qu'un automorphisme\var{\inserted{u}}{u}{insertion}{ega25_25_128-v0022} de $\overset{\hand{\vee}}{P}$ est
déterminé quand on connait \var{\xout{xxx}}{}{typewriter strike-out}{ega25_25_128-v0023}\var{\inserted[typed]{le}}{le}{insertion}{ega25_25_128-v0024} composé\var{\inserted[typed]{de\var{\xout{\unsure{s}}}{}{typewriter strike-out}{ega25_25_128-v0026} son contragrédient $\overset{\hand{\vee}}{u}$ dans P}}{de\var{\xout{\unsure{s}}}{}{typewriter strike-out}{ega25_25_128-v0026} son contragrédient $\overset{\hand{\vee}}{u}$ dans P}{insertion}{ega25_25_128-v0025} avec f: X$\hand{\longrightarrow}$ P , - et que l'asser-%
tion analogue est vraie après tout changement de base S$\hand{\rightarrow}$Spec(k)
pour un  automorphisme u de $\overset{\hand{\vee}}{P}_S$  .  Or cela résulte aussitôt  du lemme

\mspage{\hand{128}}
